% !TEX root = main.tex


This section describes the data preprocessing, cross-validation protocols,
hyperparameter selection, and experimental setup used to evaluate the four
gesture recognition methods.
 
% ----------------------------------------------------------------
\subsection{Data preprocessing}
% ----------------------------------------------------------------
 
All gestures were loaded from the raw data files and preprocessed in a
consistent manner across methods. The preprocessing pipeline consisted of three
main steps:
 
\paragraph{Normalization.} Each gesture trajectory was normalized using
standardization (z-score normalization) computed exclusively on the training
set to prevent data leakage. For each fold, the mean $\boldsymbol{\mu}$ and
standard deviation $\boldsymbol{\sigma}$ were calculated from all training
gesture points, and then applied uniformly to both training and test sets:
\begin{equation}
    \mathbf{r}^{\text{norm}}_i = \frac{\mathbf{r}_i - \boldsymbol{\mu}}{\boldsymbol{\sigma}} .
\end{equation}
This removes scale and location artifacts, ensuring that all methods operate on
gesture trajectories with zero mean and unit variance.
 
\paragraph{Dimensionality reduction.} To assess the effect of
dimensionality reduction, we optionally applied Principal Component Analysis
(PCA) to the training set before classification. PCA was fit only on the
training data and the resulting transformation was applied to both training and
test sets. We tested PCA with $n_{\text{components}} \in \{2, 3\}$ as well
as the baseline (no PCA, keeping all 3 dimensions). For gestures with
variable-length trajectories, PCA was applied point-wise, preserving the
temporal structure. PCA was not applied to the BiLSTM, as the network learns its own
feature representation directly from the resampled 3D coordinates.
 
\paragraph{Variable-length handling.} Gesture trajectories in
both domains exhibit variable length (number of
points). For methods requiring fixed-length input
(Three-Cents recognizer and BiLSTM), resampling
into a fixed number of equidistant points was
performed as part of the method-specific
preprocessing (described below). DTW and
Edit-distance handle variable-length sequences
natively without modification.
 
% ----------------------------------------------------------------
\subsection{Vector quantization for edit-distance}
% ----------------------------------------------------------------
 
The Edit-distance baseline requires discrete symbolic sequences rather than
continuous coordinates. We converted each trajectory into a sequence of cluster
labels using $k$-means clustering.
 
\paragraph{Clustering procedure.} For each fold, $k$-means clustering was fit on
all training gesture points (pooled across all gestures and users) with a fixed
random seed (42) for reproducibility. We tested cluster counts
$K \in \{15, 20, 25, 30, 35, 40\}$ to identify the optimal
discretization granularity. After fitting, each point in both training and test
sets was assigned to the nearest cluster centroid, transforming each gesture
into a string of cluster labels (e.g., $A \to A \to B \to C \to \dots$).
 
\paragraph{Compression.} We also tested an optional compression step that
removed consecutive duplicate symbols (e.g., $A \to A \to B$ becomes
$A \to B$), reducing edit-distance sensitivity to repetitive gestures. This
was tested as a binary option: compression enabled or disabled.
 
% ----------------------------------------------------------------
\subsection{Cross-validation and hyperparameter selection}
% ----------------------------------------------------------------
 
All methods were evaluated using two complementary cross-validation protocols to
assess both user-independence and generalization within users.
 
\paragraph{User-independent evaluation.} We performed leave-one-user-out
cross-validation: for each of the 10 subjects in turn, we held out all gestures
from that subject as the test set and trained on all gestures from the
remaining 9 subjects. This was repeated 10 times (once per subject), yielding
10 accuracy measurements per method/domain/hyperparameter combination.
 
\paragraph{User-dependent evaluation.} We performed leave-one-gesture-sample-out
cross-validation: for each user and each gesture type, we iteratively held out
one sample (repetition) as the test set, reserving the other 9 repetitions from
that user for training. This was repeated 10 times (once per repetition index),
again yielding 10 accuracy measurements per configuration.
 
\paragraph{Inner validation split for hyperparameter selection.}
To avoid any information leakage from the outer test fold into the
hyperparameter search, we introduced a nested validation step. Within each outer
fold, the training portion is further split randomly into an \emph{inner
training} set (80\,\%) and an \emph{inner validation} set (20\,\%), using a
fixed random seed (42) for reproducibility:
\begin{equation}
    \mathcal{D}_{\text{train}}^{\text{outer}}
    \;\xrightarrow{\;80 / 20\;}\;
    \mathcal{D}_{\text{inner-train}} \;\cup\; \mathcal{D}_{\text{inner-val}} .
\end{equation}
All preprocessing statistics (mean, standard deviation, $k$-means centroids,
PCA components) and the method parameters are fitted \emph{exclusively} on
$\mathcal{D}_{\text{inner-train}}$. Candidate hyperparameter configurations are
then scored on $\mathcal{D}_{\text{inner-val}}$. The outer test fold is thus
\emph{never} seen during hyperparameter selection, ensuring an unbiased
performance estimate. Once the best configuration is identified, the full outer
training set $\mathcal{D}_{\text{train}}^{\text{outer}}$ is used for the final
model fit before evaluating on the held-out test fold.

\paragraph{Hyperparameter grid search.} Within each outer fold, we performed a
grid search over the inner validation set, covering hyperparameters specific to
each method:

\begin{itemize}\itemsep 0em
    \item \textbf{Edit-distance:} $K \in \{15, 20, 25, 30, 35, 40\}$
    (cluster count), compression $\in \{\text{True, False}\}$,
    $k \in \{1, 3, 5, 7\}$ for $k$-NN, and
    $n_{\text{components}} \in \{\text{no PCA}, 2, 3\}$.

    \item \textbf{DTW:} $k \in \{1, 3, 5, 7\}$ for $k$-NN and
    $n_{\text{components}} \in \{\text{no PCA}, 2, 3\}$;
    no method-specific hyperparameters.

    \item \textbf{Three-Cents:} $n_{\text{points}} \in \{16, 32, 64\}$
    equidistant points and
    $n_{\text{components}} \in \{\text{no PCA}, 2, 3\}$.

    \item \textbf{BiLSTM:} the number of resampled time steps
    $T \in \{16, 32, 64\}$, the number of hidden units per direction
    $h \in \{16, 32, 64\}$, and the dropout rate
    $p \in \{0.1, 0.2, 0.3, 0.5\}$ were tuned jointly by grid search.
    The following architectural parameters were fixed throughout:
    batch size $= 16$, maximum epochs $= 50$ with early stopping
    (patience $= 5$, monitoring validation loss). No PCA was applied.
\end{itemize}

For each outer fold, we computed inner-validation accuracy for all
hyperparameter combinations. The best configuration per
(method, domain, cv\_mode) tuple was selected by maximizing mean inner-validation
accuracy across folds, and then applied to the corresponding outer test fold.
 
% ----------------------------------------------------------------
\subsection{Classification and distance metrics}
% ----------------------------------------------------------------
 
\paragraph{$k$-Nearest neighbors (\textit{k}-NN).} DTW and Edit-distance both use
$k$-NN for classification. For each test gesture, we computed distances to all
training gestures, sorted the results, and assigned the test gesture the
majority-vote label among its $k$ nearest neighbors. We tested $k \in \{1, 3,
5, 7\}$ to balance between local (small $k$) and global (large $k$)
neighborhood effects.
 
\paragraph{Distance computation.} For DTW, we used the accumulated dynamic
programming distance $D(N, M)$ as the distance metric. For Edit-distance, we
used the Levenshtein edit distance computed via dynamic programming. For
Three-Cents, we used the average point-to-point Euclidean distance.
 
\paragraph{Three-Cents 1-NN.} The Three-Cents recognizer uses 1-nearest-neighbor
by design; the $k$ parameter was fixed at 1 for this method.
 
% ----------------------------------------------------------------
\subsection{Implementation details and reproducibility}
% ----------------------------------------------------------------
 
All methods were implemented in Python 3 using NumPy for numerical computation.
DTW and Edit-distance distance computations were implemented from scratch
(without external gesture recognition libraries) as required by the project
guidelines. PCA and $k$-means clustering were sourced from scikit-learn.
 
\paragraph{Random seeds.} The $k$-means clustering algorithm was initialized
with random seed 42, and 10 different centroid initializations were tested per
fold (n\_init=10) to ensure robust cluster center identification.
 
\paragraph{Data leakage prevention.} Critical to the experimental integrity:
normalization statistics (mean, std) and clustering centroids were learned
\emph{only} from training data and applied to test data, preventing information
leakage from test to train.
 
\paragraph{Metrics.} Classification accuracy was computed as the proportion of
correct predictions:
\begin{equation}
    \text{Accuracy} = \frac{\text{Number of correct predictions}}{\text{Total number of test samples}} .
\end{equation}
For each configuration, we report mean accuracy and standard deviation across
the 10 cross-validation folds.

% ----------------------------------------------------------------
\subsection{Ablation study}
% ----------------------------------------------------------------
 
To understand the individual contribution of each preprocessing decision to final classifier performance, we conduct a systematic ablation study. In this work, two preprocessing choices are non-trivial and methodologically independent: z-score normalization (Section~4.1) and dimensionality reduction via PCA ($n_{\text{components}} \in \{2, 3\}$ or none). Rather than reporting results only for the best-performing configuration found by grid search, the ablation study evaluates all combinations of these two choices across all four classifiers, both domains, and both cross-validation protocols (80 configurations in total).

This serves three purposes. First, it allows us to assess whether the performance gains reported in Section 5.3 are genuinely driven by the chosen preprocessing or whether they would hold under alternative settings. Second, it reveals interactions between preprocessing steps that a single best-configuration table would obscure. For instance, whether PCA is beneficial in isolation or only when preceded by normalization. Third, it provides a diagnostic on the robustness of each method: a classifier whose accuracy varies widely across preprocessing choices is more sensitive to implementation decisions and harder to deploy reliably, whereas a classifier that performs consistently across configurations can be trusted without extensive tuning.


